% 简单线性回归的投影几何（示意，非按数据比例）。
% 平面表示设计矩阵 X=[1 x] 的列空间（维数2）；竖直方向代表其正交补（维数 n-2，图中只画一维）。
% 观测 y = Xβ + ε；拟合值 ŷ = Py 落在平面内，残差 e = y - ŷ = (I-P)ε 垂直于平面。
\begin{tikzpicture}[line width=.6pt,font=\fontsize{8}{10}\selectfont,>=Stealth,x={(10mm,0mm)},y={(5mm,3.4mm)},z={(0mm,10mm)}]
 % 列空间平面
 \fill[FigureBlueFill] (-0.4,-0.6,0) -- (5.4,-0.6,0) -- (5.4,3.4,0) -- (-0.4,3.4,0) -- cycle;
 \draw[FigureBlue,line width=.6pt] (-0.4,-0.6,0) -- (5.4,-0.6,0) -- (5.4,3.4,0) -- (-0.4,3.4,0) -- cycle;
 \node[FigureBlue,anchor=west] at (5.5,-0.4,0) {列空间$\operatorname{span}\{\bm 1,\bm x\}$};
 \node[FigureBlue,anchor=west] at (5.5,-0.4,-0.45) {维数$2$};
 % 原点与向量
 \coordinate (O) at (0,0,0);
 \coordinate (M) at (1.9,2.4,0);   % Xβ
 \coordinate (H) at (3.4,1.0,0);   % ŷ
 \coordinate (Y) at (3.4,1.0,2.4); % y
 \fill[FigureInk] (O) circle[radius=.7mm] node[below left,FigureInk] {$\bm 0$};
 \draw[->,FigureStroke,line width=.9pt] (O) -- (Y) node[above,FigureInk] {$\bm y$};
 \draw[->,FigureBlue,line width=.9pt] (O) -- (H) node[below right,FigureInk] {$\hat{\bm y}=\bm P\bm y$};
 \draw[->,FigureGreen,line width=.9pt] (O) -- (M);
 \fill[FigureGreen] (M) circle[radius=.7mm] node[above left,FigureGreen] {$\bm X\bm\beta$};
 % 残差（正交补方向）与噪声分解
 \draw[->,FigureRed,line width=1.1pt] (H) -- (Y) node[midway,right,FigureRed] {$\bm e=(\bm I-\bm P)\bm\varepsilon$};
 \draw[->,FigurePurple,line width=.75pt,dashed] (M) -- (Y) node[midway,left,FigurePurple] {$\bm\varepsilon$};
 \draw[->,FigurePurple,line width=.75pt,dashed] (M) -- (H) node[pos=.5,below left=-.5mm,FigurePurple] {$\bm P\bm\varepsilon$};
 % 直角标记
 \draw[FigureInk,line width=.5pt] ($(H)+(0,0,0.28)$) -- ($(H)+(-0.28,0,0.28)$) -- ($(H)+(-0.28,0,0)$);
 \node[FigureMuted,anchor=west,font=\fontsize{7.5}{9}\selectfont] at (3.6,1.0,1.85) {正交补，维数$n-2$};
\end{tikzpicture}
